\subsection{若尔当分解}

定义8. --- 设 E 是交换域上的有限维向量空间，且 u 是 E 的一个自同态。则 u 的一个若尔当分解是一对 \((u_{s}, u_{n})\)，其中 \(u_{s}\) 是 E 的半单自同态，\(u_{n}\) 是 E 的幂零自同态，使得 \(u_{s}u_{n} = u_{n}u_{s}\) 且 \(u = u_{s} + u_{n}\)。

定理1. --- 设 E 是交换域 K 上的有限维向量空间，且设 u 是 E 的一个自同态。则 u 具有若尔当分解 \((u_{s}, u_{n})\) 当且仅当 u 的特征值在 K 上是可分的。此外，该分解是唯一的，u 与 \(u_{s}\) 的特征多项式相同，并且存在没有常数项的多项式 P 和 \(Q \in K[X]\)，使得 \(u_{s} = P(u)\) 且 \(u_{n} = Q(u)\)。

\begin{enumerate}
\def\labelenumi{\Alph{enumi})}
\tightlist
\item
  让我们首先证明以下特殊情况：
\end{enumerate}

引理4. --- 设 E 是交换域 K 上的有限维向量空间，且设 u 是 E 的一个可三角化的自同态。则存在唯一的与 u 交换的可对角化自同态 v，使得 u - v 是幂零的。此外，在这些条件下，u 和 v 的特征多项式相同，并且存在多项式 \(P \in K[X]\) 使得 \(v = P(u)\)。

设 \(v\) 是 \(E\) 的一个可对角化自同态，使得 \(uv = vu\) 且 \(v - u\) 是幂零的；设 \(\alpha\) 是 \(v\) 的一个特征值，并令 \(V_{\alpha}\) 为对应的特征子空间。由引理 3（VII，第 41 页），\(V_{\alpha}\) 在 \(u\) 下封闭，且 \(u - \alpha\) 在 \(V_{\alpha}\) 上的限制也是 \(u - v\) 的限制，因而是幂零的；所以 \(V_{\alpha}\) 包含在子空间 \(M_{\alpha}\) 中，后者由那些 \(x \in E\)（被 \(u - \alpha\) 的某个幂零化的）组成。由于 \(E\) 是诸 \(V_{\alpha}\) 的直和，也是诸 \(M_{\alpha}\) 的直和（VII，第 34 页，命题 7），这表明 \(V_{\alpha} = M_{\alpha}\)（对所有 \(\alpha\)）。由命题 9 的推论（VII，第 36 页）可得 \(\chi_u = \chi_v\)；同样可得 \(v\) 由 \(u\) 唯一确定；它在每个 \(M_{\alpha}\) 上的限制是由 \(\alpha\) 定义的位似变换。

反之，让我们按上述条件定义 \(v\)；显然 \(v\) 是可对角化的，且 \(u - v\) 是幂零的。由 VII，第 34 页，命题 7，存在多项式 \(q_{\alpha}\)，使得对所有 \(x \in E\)，\(x\) 在 \(M_{\alpha}\) 中的分量是 \(q_{\alpha}(u) \cdot x\)。于是 \(v = \sum_{\alpha} \alpha q_{\alpha}(u)\)；这蕴含 \(u\) 与 \(v\) 可交换，从而完成证明
% semantic-fragments: legacy-input-seam
\relax{}
\textbf{B) 现在让我们回到定理1的证明。}\par

首先，假设 \(u\) 可以写成形式 \(s + n\) ，其中 \(s\) 是绝对半单的，而 \(n\) 是幂零的，并且 \(s\) 与 \(n\) 可交换。设 \(\Omega\) 为 \(K\) 的一个代数闭包；则 \(u_{(\Omega)} = s_{(\Omega)} + n_{(\Omega)}\) ，其中 \(s_{(\Omega)}\) 是可对角化的，而 \(n_{(\Omega)}\) 是幂零的，并且 \(s_{(\Omega)}\) 与 \(n_{(\Omega)}\) 可交换；由引理4可知，\(s_{(\Omega)}\) 从而 \(s\) 也一样，是唯一的，多项式 \(\chi_{u_{(\Omega)}}\) 与 \(\chi_{s_{(\Omega)}}\) 在 \(\Omega[X]\) 中一致，因此多项式 \(\chi_u\) 与 \(\chi_s\) 也一致，并且 \(s\) 可以表示为 \(u\) 的、系数属于 \(\Omega\) 的多项式。这首先表明，\(u\) 的特征值与 \(s\) 的特征值相同，从而在 \(K\) 上可分（VII，第42页，命题15）；此外，由于 \(s\) 是一个 \(\Omega\) -线性组合，由 \(u\) 的幂组成，它也是同样这些幂的 \(K\) -线性组合（II，第～311页，命题～19），故存在多项式 \(P \in K[X]\) 使得 \(s = P(u)\) ，从而 \(n = Q(u)\) ，其中 \(Q(X) = X - P(X)\) 。现在让我们证明 \(Q\) （从而 \(P\) ）可以取得没有常数项。若 \(u\) 可逆，则其特征多项式有非零常数项，且凯莱-哈密顿定理（III，第～541页，命题～20）表明 1 可以表示为 \(u\) 的一个没有常数项的多项式，故断言在此情形成立。若 \(u\) 不可逆，则其核 \(W\) 非零且在 \(n\) 下封闭（VII，第41页，引理3）；由于 \(n\) 在 \(W\) 上的限制是幂零的，故存在向量 \(x \neq 0\) 属于 \(W\) ，使得 \(u(x) = n(x) = 0\) ，这表明 \(Q\) 不能有常数项。

反之，假设 \(u\) 的特征值在 \(K\) 上可分，并设 \(L\) 是包含这些特征值的、\(K\) 的有限伽罗瓦扩张。根据引理4，我们可以写出 \(u_{(L)} = v + w\) ，其中 \(v\) 是可对角化的，而 \(w\) 是幂零的，并且 \(vw = wv\) 。设 \(B\) 为 \(E\) 的一组基，设 \(B'\) 为 \(L \otimes_K E\) 的对应基，并设 \(U, V, W\) 为 \(u_{(L)}, v, w\) 关于 \(B'\) 的矩阵；注意， \(U\) 也是 \(u\) 关于 \(B\) 的矩阵，故其元素属于 \(K\) 。对每个 \(K\) -自同构 \(\sigma\) （即 \(L\) 的自同构），以及每个矩阵 \(A\) ，其元素属于 \(L\) ，令 \(A^\sigma\) 表示把 \(\sigma\) 作用于 \(A\) 的元素所得到的矩阵。设 \(\sigma\) 是一个 \(K\) -自同构，即 \(L\) 的自同构；则 \(U = U^\sigma = (V + W)^\sigma = V^\sigma + W^\sigma\) , \(V^\sigma W^\sigma = (VW)^\sigma = (WV)^\sigma = W^\sigma V^\sigma\) ；由于 \(V^\sigma\) 是某个可对角化自同态的矩阵，而 \(W^\sigma\) 是幂零的，由引理4可知 \(V^\sigma = V\) 且 \(W^\sigma = W\) 。由于这对所有 \(\sigma\) 都成立，V 和 W 的元素属于 K；若 \(u_{s}\) 与 \(u_{n}\) 是 E 的、关于基 B 的矩阵分别为 V 和 W 的自同态，则 \((u_{s})_{(L)} = v\) 且 \((u_{n})_{(L)} = w\) 。由此可知， \(u_{s}\) 是绝对半单的， \(u_{n}\) 是幂零的， \(u_{s}\) 与 \(u_{n}\) 可交换，并且 \(u = u_{s} + u_{n}\) 。这就完成了证明。

每当一个自同态 \(f\) 具有若尔当分解时，我们将其记为 \((f_s, f_n)\) ，并且自同态 \(f_s\) 与 \(f_n\) 分别称为 \(f\) 的绝对半单分量和幂零分量。当 \(K\) 是完全域时，每个自同态都有若尔当分解；在这种情况下，绝对半单自同态与半单自同态之间也没有区别，我们有时将“绝对半单分量”称为“半单分量”。

推论1. --- 设u有若尔当分解，且L是K的一个扩域。则 \(u_{(L)}\) 具有若尔当分解，其中 \((u_{(L)})_{s} = (u_{s})_{(L)}\) 且 \((u_{(L)})_{n} = (u_{n})_{(L)}\) .

推论2. --- 设u有若尔当分解。则E的每个与u可交换的自同态也与 \(u_{s}\) 和 \(u_{n}\) 可交换.

推论3. --- 设u和v是E的两个可交换的自同态，且它们具有若尔当分解。

\begin{enumerate}
\def\labelenumi{\alph{enumi})}
\item
  这些自同态 \(u, v, u_s, v_s, u_n, v_n\) 均交换。
\item
  这些自同态 \(u + v\) 与 \(uv\) 具有若尔当分解，且
\end{enumerate}

\[
\begin{array}{c} (u + v) _ {s} = u _ {s} + v _ {s}, \quad (u + v) _ {n} = u _ {n} + v _ {n}, \quad (u v) _ {s} = u _ {s} v _ {s}, \\ (u v) _ {n} = u _ {s} v _ {n} + u _ {n} v _ {s} + u _ {n} v _ {n}. \end{array}
\]

a) 部分由推论2得出。要证明 b) 部分，只需注意到 \(u_{s} + v_{s}\) 与 \(u_{s}v_{s}\) 是绝对半单的（VII，第43页，推论），且 \(u_{n} + v_{n}\) 与 \(u_{s}v_{n} + u_{n}v_{s} + u_{n}v_{n}\) 是幂零的（作为可交换的幂零自同态之和）。

推论4. --- 设u有若尔当分解，并设 R 是 K{[}X{]} 中的多项式。则自同态 \(\mathrm{R}(u)\) 有若尔当分解，且 \(\mathrm{R}(u)_{s} = \mathrm{R}(u_{s})\) .

注记. --- 1) 我们有 \(\det(u_s) = \det(u)\) 和 \(\operatorname{Tr}(u_s) = \operatorname{Tr}(u)\) .

\begin{enumerate}
\def\labelenumi{\arabic{enumi})}
\setcounter{enumi}{1}
\tightlist
\item
  u 可三角化的一个充分必要条件是：u 有一个若尔当分解，且 \(u_{s}\) 可对角化。则存在 E 的一组基，使得 u 关于该基的矩阵是下三角的，而 \(u_{s}\) 关于该基的矩阵是对角的，且与 u 的矩阵具有相同的对角线（参见下面的引理4和命题19）。
\end{enumerate}

但要注意，如果u关于某组基的矩阵是三角的，一般不能推出 \(u_{s}\) 关于同一组基的矩阵是对角的。

\begin{enumerate}
\def\labelenumi{\arabic{enumi})}
\setcounter{enumi}{2}
\tightlist
\item
  方阵的若尔当分解的概念可以类似地定义。例如，对于若尔当矩阵 \(U_{m,\alpha}\) 我们有
\end{enumerate}

\[
(U _ {m, \alpha}) _ {s} = \alpha . I _ {m}, (U _ {m, \alpha}) _ {n} = U _ {m, 0}.
\]

\begin{enumerate}
\def\labelenumi{\arabic{enumi})}
\setcounter{enumi}{3}
\tightlist
\item
  如果 \(u\) 是半单的但不是绝对半单的，则它没有若尔当分解。
\end{enumerate}

交换域上向量空间 V 的自同态 u 称为幂单的，如果自同态 \(u - Id_{V}\) 是幂零的，即存在整数 r 使得 \((u - Id_{V})^{r} = 0\) ；则 u 是 V 的自同构，因为若 \(u = Id_{V} - n\) ，其中 \(n^{r} = 0\) ，则

\[
\left(\operatorname{Id} _ {\mathrm{V}} + n + \dots + n ^ {r - 1}\right) u = u \left(\operatorname{Id} _ {\mathrm{V}} + n + \dots + n ^ {r - 1}\right) = \operatorname{Id} _ {\mathrm{V}}.
\]

若 V 的维数为 m，则 u 是幂单的当且仅当 \(\chi_{u}(X)=(X-1)^{m}\) （VII，第33页，命题5的推论3）。

命题17. --- 设E是交换域上的有限维向量空间，并设f是E的自同态。则以下条件等价：

\begin{enumerate}
\def\labelenumi{(\roman{enumi})}
\item
  \(f\) 有若尔当分解且是自同构；
\item
  \(f\) 有若尔当分解且 \(f_{s}\) 是自同构；
\item
  存在一个绝对半单的自同构 \(a\) 作用于 \(E\) ，和一个幂单的自同态 \(u\) 作用于 \(E\) ，使得 \(f = ua = au\) .
\end{enumerate}

此外，在这些条件下，在(iii)的记号中，我们必须有 \(a = f_{s}\) 和 \(u = 1 + f_{s}^{-1}f_{n}\) .

\begin{enumerate}
\def\labelenumi{(\roman{enumi})}
\item
  \(\Rightarrow\) （ii）：这由注记1得出。
\item
  \(\Rightarrow\) （iii）：取 \(a = f_{s}\) 和 \(u = 1 + f_{s}^{-1}f_{n}\) ；则 \(f = ua = au\) ，而 \(a\) 是一个绝对半单自同构，且 \(u\) 是幂单的。
\item
  \(\Rightarrow\) (i)：在(iii)的记号下，取 \(n = a(u - 1) = (u - 1)a\) ，则 \(an = na\) 且 \(f = a + n\) ，且 \(n\) 是幂零的。由此可知， \((a, n)\) 是 \(f\) 的若尔当分解。这蕴含了(i)以及关系式 \(a = f_s\) 和 \(u = 1 + f_s^{-1}f_n\) .
\end{enumerate}

令 \(f_{u}=f_{s}^{-1}f=ff_{s}^{-1}=1+f_{s}^{-1}f_{n}\) ，并称它为 f 的幂单分量。数对 \((f_{s},f_{u})\) 通常称为自同构 f 的乘法若尔当分解。

命题18. --- 设E是交换域K上的有限维向量空间，设u是E的一个自同态，并设E'是E中在u下封闭的子空间。设u'（分别地，u''）是由u诱导的E'（分别地，E/E'）的自同态。则 \(\chi_{u} = \chi_{u'} \cdot \chi_{u''}\) .

为使 u 具有若尔当分解，必要且充分条件是 \(u'\) 和 \(u''\) 具有；此外，若此成立，则 \(u'\) 和 \(u''\) 的绝对半单（分别地，幂零）分量是 \(E'\) 和 \(E/E'\) 的自同态，它们由 u 的绝对半单（分别地，幂零）分量诱导。

设 B 为 E 的一个基，它包含一个基 \(B'\) ，即 \(E'\) 的一个基，并设 \(B''\) 是

\(E'' = E/E'\) 的基，即 B - B' 的像。设 U, \(U'\) , \(U''\) 分别是 u, \(u'\) , \(u''\) 关于 B, \(B'\) , \(B''\) 的矩阵。则 U 具有形式

\[
\left( \begin{array}{c c} U ^ {\prime} & Z \\ 0 & U ^ {\prime \prime} \end{array} \right)
\]

且 \(\chi_{u} = \chi_{U} = \chi_{U'} \chi_{U''} = \chi_{u'} \chi_{u''}\) （参见 III，第 533 页，习题 2）。我们推断 u 的特征值集合是 \(u'\) 和 \(u''\) 的特征值集合的并。若 \(u'\) 和 \(u''\) 具有若尔当分解，则 \(u'\) 和 \(u''\) 的特征值在 K 上可分，因此 u 的特征值也是可分的，且 u 具有若尔当分解（VII，第43页，定理1）。反之，若 u 具有若尔当分解 (s, n)，则 s 和 n 保持 \(E'\) 不变，因为它们是 u 的多项式；设 \(s', n', s'', n''\) 表示 \(E', E', E'', E''\) 的、分别由 s、n、s、n 诱导的自同态。则 \(s'\) 和 \(s''\) 是绝对半单的，因为它们的极小多项式整除 s 的极小多项式（VII，第42页，命题15）；此外 \(n'\) 和 \(n''\) 是幂零的。最后 \(u' = s' + n', u'' = s'' + n'', s'n' = n's',\) 且 \(s''n'' = n''s''\) ，这就完成了证明。

命题19. --- 设E是交换域K上的有限维向量空间，且设F是由E的两两可交换的可三角化自同态组成的集合。则存在E的一组基，使得F中每个元素u关于这组基的矩阵都是下三角的，且 \(u_{s}\) 关于这组基的矩阵是对角的，其对角元素与u的相同。

根据VII第45页推论3，集合 \(\mathfrak{F}_s\) ——即 \(\mathfrak{F}\) 的元素的绝对半单分量组成的集合——由彼此交换的可对角化元素组成，故该集合可对角化（VII，第41页，命题13）；集合 \(\mathfrak{F}_n\) ——即 \(\mathfrak{F}\) 的元素的幂零分量组成的集合——由两两交换的幂零元素组成，且 \(\mathfrak{F}_n\) 的每个元素与 \(\mathfrak{F}_s\) 的每个元素交换。按照命题13（第七章，第41页）证明中的论证方式，我们看到存在将E分解为子空间的直和 \(E_i\) ，这些子空间在 \(\mathfrak{F}_s\) 和 \(\mathfrak{F}_n\) 下不变，并且使得 \(\mathfrak{F}_s\) 的每个元素在每一个 \(E_i\) 上的限制都是一个位似变换。将 E 依次替换为每个 \(E_i\) ，我们可以假设 \(\mathfrak{F}_s\) 的元素都是位似变换；只需证明存在 E 的一组基，使得 \(\mathfrak{F}_n\) 的元素在这组基下由对角线为零的下三角矩阵表示；这样我们就归结为 \(\mathfrak{F}\) 由幂零元素组成的情形。

现在假设 \(E \neq 0\) ，并设 F 为 E 的一个非零子空间，在 F 下不变，且维数最小。则对每个 \(u \in F\) ，u 在 F 上的限制的核非零且在 F 下不变（VII，第 41 页，引理 3）；由 F 的选取，u 在 F 上的限制从而对所有 \(u \in F\) 都为零。设 \(x \in F\) , \(x \neq 0\) ；则 \(u(x) = 0\) 对所有 \(u \in F\) 成立；通过对E的维数进行归纳论证，我们可以假设存在一个基 \((\bar{e}_{1}, \ldots, \bar{e}_{n-1})\) （即商 \(E' = E/Kx\) 的基），使得，对于所有 \(u \in F\) ，由 u 诱导的自同态 \(\bar{u}\) （ \(E'\) 上的）关于该基的矩阵是下三角的且对角线为零；如果 \(e_{i} \in E\) 投影到 \(\bar{e}_{i}\) 上，对于 i = 1， \ldots，n - 1，则基 \((e_{1}, \ldots, e_{n-1}, x)\) 满足所需条件。

